% ECONOMICS_PROBLEM: {"id": "I26-412", "title": "Returns to expanding college access", "jel_code": "I26", "source_id": "OP-0454"}
% !TeX program = xelatex
\documentclass[11pt,a4paper]{article}
\usepackage[margin=23mm]{geometry}
\usepackage{fontspec}
\setmainfont{Latin Modern Roman}
\usepackage{amsmath,amssymb,mathtools}
\usepackage{xcolor,enumitem,booktabs,longtable,array,xurl,hyperref}
\definecolor{muted}{HTML}{53636C}
\hypersetup{colorlinks=true,urlcolor=blue,linkcolor=blue}
\setlength{\parindent}{0pt}
\setlength{\parskip}{5pt}
\newcommand{\R}{\mathbb R}
\newcommand{\E}{\mathbb E}
\newcommand{\N}{\mathbb N}
\newcommand{\ind}{\mathbf 1}
\newcommand{\Var}{\operatorname{Var}}
\newcommand{\Cov}{\operatorname{Cov}}
\newcommand{\supp}{\operatorname{supp}}
\DeclareMathOperator*{\argmin}{arg\,min}
\DeclareMathOperator*{\argmax}{arg\,max}
\newcommand{\entrymeta}[3]{{\sffamily\small\color{muted}\textbf{\texttt{#1}}\quad|\quad #2\par #3\par}}
\newcommand{\Problemref}[1]{\texttt{#1}}
\newcommand{\JELref}[2]{\texttt{#2} (JEL \texttt{#1})}
\begin{document}
\section*{Returns to expanding college access}
\textbf{Problem ID:} \texttt{I26-412}\quad \textbf{JEL:} \texttt{I26}\par
% BEGIN ECONOMICS STATEMENT
\label{problem:OP-0454}
{\sffamily\small\color{muted}Original subject group: Education, families and demography\par}
\label{definitions:problem:30007736}
\textbf{Audit classification:} Background; proposed extension. Distinguished capacity-wide applicant net returns from the already estimated local admission margin and national welfare.
\subsection*{Definitions and precise research target}
\textit{Editorial formalization.} Consider applicants near the admission threshold of one public-university system. Let \(k\in\{0,1,\ldots,\bar k\}\) be an administratively feasible integer capacity expansion, where \(\bar k\) is a declared maximum measured in additional seats, and let \(E_i(k)\) be applicant \(i\)'s enrollment destination after admissions, financial aid, and peer composition adjust. Define discounted lifetime earnings \(Y_i(k)\), education resource cost \(C_i(k)\), and a prespecified monetary nonpecuniary benefit \(N_i(k)\). For the baseline applicant cohort, define
\[V(k)=E\!\left[\sum_i\{Y_i(k)+N_i(k)-C_i(k)\}\right]-K(k),\]
where \(K(k)\) is additional infrastructure cost not already included in \(C_i(k)\). Determine the discrete marginal net value \(V(k+1)-V(k)\) for \(k<\bar k\) and distinguish it from a local admissions-cutoff earnings effect. The identification strategy must address alternative colleges, displacement of other applicants, completion, peer effects, and declining instructional resources per pupil. Lifetime quantities beyond observed follow-up require explicit extrapolation and uncertainty intervals. Report separately outcomes for marginal admits and changes for existing students.

Fix the full baseline applicant set, including existing admits, rejected applicants and substitutes at other institutions. \(k\) adds exactly \(k\) seats under a stated admissions/aid rule, with capacity and peer adjustments declared. \(Y_i(k)=\sum_{t=1}^{T}\beta^ty_{it}(k)+\operatorname{terminal}_i(k)\); the terminal term represents a stated lifetime extrapolation, not observed earnings. \(N_i(k)\) is a monetary nonpecuniary benefit under a common reference convention, and \(C_i,K\) partition direct educational and infrastructure opportunity costs. Earnings are a private-income measure, not automatically national output: specify whether price/general-equilibrium wage changes are excluded or incorporate all taxpayers and affected workers for a social welfare extension. The stated \(V\) is an applicant-cohort net-return objective, rather than a universal utilitarian welfare functional. All expectations use a stated baseline population and probability law. Identification means every admissible data-generating model with the same observed-data law yields the same displayed target; otherwise report the identified set. Exact equations, horizons and assumptions are editorial, rather than source-given canonical conjectures.
\subsection*{Variants and assumption changes}
\begin{enumerate}
\item \textbf{Editorial alternative-college variant.} Stratify marginal admits by their no-expansion destination using explicit principal-stratum assumptions or bounds. Gains over no college and over another college can differ.
\item \textbf{Editorial national-welfare variant.} Include affected workers, taxpayers and owners with endogenous wages and baseline-money welfare. Compare with applicant earnings under the same financing rule.
\end{enumerate}
\subsection*{References and source audit}
\textbf{1.} Existing local admissions-cutoff result; capacity-wide marginal welfare remains a different editorial target. Locator: Working-paper abstract; revision December 2025. Primary indexed abstract inspected; direct page retrieval failed. URL: \url{https://www.nber.org/papers/w32296}.
\begin{enumerate}\raggedright
\item\hypertarget{definitions:source:30007736:1}{} Jack Mountjoy. 2024. \emph{Marginal Returns to Public Universities}. Working-paper abstract; revision December 2025.
\url{https://www.nber.org/papers/w32296}.
DOI: \url{https://doi.org/10.3386/w32296}.
\end{enumerate}

\subsection*{Common conventions: Source mathematical conventions}

These conventions apply to each entry unless explicitly replaced.

\textbf{Models and identification.} A primitive model \(m\) specifies all probability laws, feasible choices, information, equilibrium and measurement rules used by its target. The admissible class \(\mathcal M\) is fixed independently of outcomes. If \(P_O(m)\) is its observable probability law and \(\tau(m)\) the target, its identified set at observable law \(P\) is
\[
 \mathcal I_\tau(P)=\{\tau(m):m\in\mathcal M,\ P_O(m)=P\}.
\]
Point identification means that this set is a singleton. An empty set signals incompatibility of data and assumptions. Coordinate bounds need not describe the joint set. Bounds are sharp only if every retained target is attainable or arbitrarily approachable by compatible models. Finite bounds require explicit restrictions on unobserved quantities; unrestricted confounding, hidden wealth, extrapolation or arbitrary equilibrium selection can yield uninformative sets.

\textbf{Causal interventions.} A potential outcome \(Y_i(p)\) is the outcome under a complete policy regime \(p\), including timing, financing, information and market spillovers. The target population and units are fixed before treatment. With interference, use the complete assignment vector or a justified exposure mapping; individual no-interference notation alone cannot identify an economy-wide intervention. A difference of mean potential outcomes does not require the cross-world joint distribution. Conditioning on post-treatment migration, survival, enrollment or employment changes the estimand and needs separate justification.

\textbf{Statistical statements.} An estimator, confidence set, forecast or test specifies a sampling law, model class and loss/coverage criterion. Uniform \((1-\alpha)\)-coverage means \(\inf_{m\in\mathcal M}\Pr_m\{\tau(m)\in C_n\}\geq1-\alpha\), or an explicitly asymptotic version. The whole parameter space trivially has coverage, so informativeness requires a stated width, risk or power objective. A forecasting improvement is relative to a named benchmark, information set and evaluation population.

\textbf{Equilibrium and optimization.} An equilibrium comprises feasible plans, optimality, beliefs where needed, budget constraints and all market/resource clearing. The primitives or a stated benchmark define each correspondence \(\mathcal E(p;m)\). Nonemptiness is assumed or proved, never inferred just from compact policy sets. A selected-equilibrium target uses a declared selection rule; a robust target quantifies over all permitted equilibria. Use suprema or infima unless attainment is established. An \(\varepsilon\)-optimal policy is feasible and within \(\varepsilon\) of the finite optimum value. Report an empty feasible set separately.

\textbf{Welfare and accounting.} Normative utility, interpersonal normalization, social weights, numeraire, discounting and the use of public receipts are primitives. Behavioral utility need not coincide with the welfare criterion. Household welfare includes ownership income and rents once; adding profits again to compensating variation generally double-counts them. A monetary equivalent is defined by an explicit transfer and price convention, with monotonicity and range conditions. Average effects, mechanism contributions, transfers and welfare losses are different targets. Infinite sums require convergence and dynamic feasible plans require terminal or no-Ponzi conditions.

\textbf{Source versions.} A working-paper date, revision date and journal publication year are distinguished. A DOI for a journal version is not automatically the DOI of an earlier manuscript. Locators refer to the stated version and printed pages where specified. A metadata-only check does not verify a claim about a full-text conclusion. Every entry's source subsection narrows the provenance claim accordingly.
% END ECONOMICS STATEMENT
\end{document}
